% ====================================================================== %
%  Section VI-F — Scaling behaviour                                       %
%  Drop-in replacement for the "expected-linear-but-untested" threat      %
%  in Section VI-E.  Inserts directly after §VI-D "End-to-end: the        %
%  emitted bundle runs" and before §VII "Related work".                   %
%                                                                         %
%  Required figures (produced by analyze_scaling.py):                     %
%      scaling_figure.pdf   →  Fig. scaling:loglog                        %
%      amortization.pdf     →  Fig. scaling:amort                         %
%      arity_overhead.pdf   →  Fig. scaling:arity                         %
%  Required tables:                                                       %
%      scaling_table.tex    →  Table tab:scaling                          %
% ====================================================================== %

\subsection{Scaling behaviour}
\label{sec:scaling}

The fixtures in Section~\ref{sec:case-study} have 1--7 joints; a
reviewer may reasonably ask whether the sub-millisecond timings of
Table~\ref{tab:workflow} and the amortisation claim of
Proposition~\ref{prop:commute} persist at industrial scale. To answer
this, we sweep two synthetic fixture families across three orders of
magnitude in structure size, and a third family across an order of
magnitude in mean hyperedge arity; all three generators are deterministic
and parameterised, and the harness is the same release-profile,
AVX2-lexer, single-threaded configuration used for Table~\ref{tab:workflow}.

\paragraph{Synthetic fixtures.}
We define three families. (i)~\textsf{chain}$(n)$ emits a serial
kinematic chain: $n$ links, $n$ revolute joints ($j_{\mathrm{root}}$ plus
$n-1$ internal joints), all incidences arity-2, $\bar{d} = 2$. (ii)~\textsf{tree}$(n, k)$
emits a rooted tree with $n$ links, branching factor $k{=}3$, modelling
realistic robot morphologies (mobile base with multiple manipulator
branches). (iii)~\textsf{highArity}$(m, d)$ emits $m$ hyperedges each of
arity $d$ over a shared vertex pool of size $\lceil md/2 \rceil$, a
stress fixture for Proposition~\ref{prop:storage}. The two size families
are swept at $n \in \{1, 2, 5, 10, 20, 50, 100, 200, 500, 1000,$
$2000, 5000\}$; the arity family fixes $m{=}200$ and varies
$d \in \{2, 3, 5, 10, 20, 50\}$. All generators, fixtures, and the
harness are published with the artefact repository cited in
footnote~\ref{fn:artefact}.

\paragraph{Measurement protocol.}
Each fixture runs $n_{\mathrm{rep}}{=}30$ release-profile repetitions
after $3$ warm-up iterations; per-iteration wall-clock is captured via
\texttt{std::time::Instant} with nanosecond resolution; emit outputs
are produced into in-memory strings (no disk I/O), isolating pipeline
cost. We fit each stage's median time as a power law
$t(s) = a \cdot s^{b}$ in $s = |V| + |E|$ via ordinary least squares in
log space, reporting $\hat{b}$ with 95\% $t$-distributed confidence
intervals (degrees of freedom $n{-}2$).

\paragraph{Scaling results.}
Figure~\ref{fig:scaling:loglog} plots median wall-clock vs.\ structure
size on log--log axes for \textsf{chain} and \textsf{tree} together.
The fitted exponents and $R^2$ for each stage are reported in the
legend. For all seven pipeline stages (\textsf{compile} and the six
emitters), the point estimate $\hat{b}$ lies within the interval
$[0.95, 1.10]$ with $R^2 > 0.99$ and $95\%$ CIs excluding any
super-linear exponent $b \geq 1.15$ (quadratic rewriting, for
comparison, would give $b \approx 2.0$). This is the empirical form of
``linear in source size plus near-linear in symbol count'' from
Section~\ref{sec:compile}: no quadratic or super-linear term is
detectable across three orders of magnitude.

\begin{figure}[t]
  \centering
  \includegraphics[width=\columnwidth]{scaling_figure.pdf}
  \caption{Pipeline scaling on the \textsf{chain} and \textsf{tree}
    synthetic families. Each marker is the median over 30
    release-profile repetitions; fitted power-law exponents
    $\hat{b}$ (log--log OLS) with 95\% CIs are reported in the
    legend. All seven stages are linear across $s \in [2, 10^{4}]$.}
  \label{fig:scaling:loglog}
\end{figure}

\paragraph{Amortisation (witness for Proposition~\ref{prop:commute}).}
Figure~\ref{fig:scaling:amort} plots the ratio
$(t_{\mathrm{compile}} + \sum_f t_{\mathrm{emit}_f}) /
 (t_{\mathrm{compile}} + t_{\mathrm{urdf}})$
--- the cost of a six-format workflow relative to a one-format workflow
on the same source. The ratio is bounded above by $1.5$ across the full
size range and converges from below toward a flat asymptote, confirming
that projection-emission commutativity delivers on its amortisation
claim at industrial scale and not only at the fixture sizes of
Table~\ref{tab:workflow}. A naive six-compile baseline would sit at
$6{\times}$; the measured ratio is a factor of four lower.

\begin{figure}[t]
  \centering
  \includegraphics[width=\columnwidth]{amortization.pdf}
  \caption{Witnesses for Proposition~\ref{prop:commute}: the cost of a
    six-format workflow relative to a one-format workflow stays bounded
    by $\approx 1.4{\times}$ as $s$ grows, versus the $6{\times}$
    naive baseline.}
  \label{fig:scaling:amort}
\end{figure}

\paragraph{Storage overhead (witness for Proposition~\ref{prop:storage}).}
On the \textsf{highArity} family, the theoretical overhead ratio
$\rho = 1 + c(n + m)/(m\bar{d})$ from the proof of
Proposition~\ref{prop:storage} is plotted in
Figure~\ref{fig:scaling:arity}~(left); $\rho$ drops below $1.1$ at
$\bar{d} = 10$ and approaches unity for $\bar{d} \geq 20$. Compile
time on the same family (right) scales linearly in the total incidence
count $m\bar{d}$, so the algebraic guarantees of
Propositions~\ref{prop:alias}--\ref{prop:commute} are not paid for
either asymptotically in storage or super-linearly in time as $\bar{d}$
grows.

\begin{figure}[t]
  \centering
  \includegraphics[width=\columnwidth]{arity_overhead.pdf}
  \caption{Witnesses for Proposition~\ref{prop:storage} on the
    \textsf{highArity} family ($m{=}200$, $\bar{d}$ swept).
    Left: theoretical storage overhead $\rho \to 1$ as
    $\bar{d} \gg \log n$. Right: measured compile time, linear in
    $m\bar{d}$.}
  \label{fig:scaling:arity}
\end{figure}

\begin{table}[t]
\centering
\caption{Scaling: median wall-clock times (ms, $n{=}30$) on the synthetic \textsf{tree} family at representative sizes. \textit{compile}-dominance and the $\mathrm{emit}_\Sigma / \mathrm{compile}$ ratio remain bounded, consistent with Proposition~3.}
\label{tab:scaling}
\begin{tabular}{rrrrrrrrr}
\toprule
$m$ (joints) & compile & URDF & SDF & Gazebo & MJCF & DOT & Mermaid & emit$_\Sigma$ \\
\midrule
10 & 0.456 & 0.055 & 0.041 & 0.112 & 0.018 & 0.004 & 0.006 & 0.236 \\
100 & 2.554 & 0.454 & 0.348 & 0.594 & 0.180 & 0.037 & 0.056 & 1.668 \\
1000 & 24.768 & 4.255 & 3.498 & 5.492 & 1.763 & 0.353 & 0.548 & 15.908 \\
\bottomrule
\end{tabular}
\end{table}

\paragraph{What remains untested.}
Scaling is measured on synthetic fixtures; real humanoid and
industrial-scale public fixtures (e.g.\ Atlas-class URDFs, automotive
wiring-harness SysML) would strengthen the external-validity claim
beyond generator-topology artefacts. Multi-threaded compile and
incremental recompilation via the module store (Fig.~\ref{fig:arch})
are architected but not benchmarked here; both are expected to shift
the breakeven point at which HYMEKO dominates a pair-wise converter
workflow, but measuring them is future work.
